% SGA 7, tome I — front matter source transcription layer.
% UTF-8 without BOM; LF line endings. Fragment only; no standalone build claimed.
% Controlling scan: PDF pages 1--12 / zero-based indices 0--11.
% Scan SHA-256: 9CD40FF06EB1E488AF385A56899D4F492492A06A1E2E3C0ED6876B82E3E3603F.
% Editorial gap and obstruction notices are not source text.

\providecommand{\SGASevenSourceGap}[1]{\textit{[lacune visuelle : #1]}}
\providecommand{\SGASevenSourceObstruction}[1]{\textit{[masqué par #1]}}

% SGA7I-FRONT-P000 — controlling PDF page 1 — cover.
\begin{center}
Lecture Notes in Mathematics

A collection of informal reports and seminars

Edited by A. Dold, Heidelberg and B. Eckmann, Zürich

\bigskip
288

\bigskip
Séminaire de Géométrie Algébrique du Bois-Marie

1967--1969

SGA 7 I

Dirigé par A. Grothendieck

avec la collaboration de M. Raynaud et D. S. Rim

\bigskip
Groupes de Monodromie

en Géométrie Algébrique

\bigskip
Springer-Verlag

Berlin : Heidelberg : New York
\end{center}
\clearpage

% SGA7I-FRONT-P001 — controlling PDF page 2 — publisher endpaper.
% The Nottingham bookplate covers the centre of the publisher overview.
% Visible fragments on both sides of the obstruction are kept without reconstructing the hidden text.
\noindent
This series aims to report new d\SGASevenSourceObstruction{l'ex-libris de bibliothèque}s - quickly, informally\par
and at a high level. The type of m\SGASevenSourceObstruction{l'ex-libris de bibliothèque}\par
1. Preliminary drafts of original p\SGASevenSourceObstruction{l'ex-libris de bibliothèque}\par
2. Lectures on a new field, or pre\SGASevenSourceObstruction{l'ex-libris de bibliothèque}\par
3. Seminar work-outs\par
4. Reports of meetings, provide\SGASevenSourceObstruction{l'ex-libris de bibliothèque}\par
\hspace*{1em}a) of exceptional interest or \SGASevenSourceObstruction{l'ex-libris de bibliothèque}\par
\hspace*{1em}b) devoted to a single topic.\par
Texts which are out of print but s\SGASevenSourceObstruction{l'ex-libris de bibliothèque}thin these categories.\par
The timeliness of a manuscript i\SGASevenSourceObstruction{l'ex-libris de bibliothèque}finished or tentative.\par
Thus, in some instances, proofs \SGASevenSourceObstruction{l'ex-libris de bibliothèque}ch have been or will\par
later be published elsewhere.\par
Manuscripts should comprise not \SGASevenSourceObstruction{l'ex-libris de bibliothèque}\par
Publication of \textit{Lecture Notes} is in\SGASevenSourceObstruction{l'ex-libris de bibliothèque}tical community, in\par
that a commercial publisher, Spri\SGASevenSourceObstruction{l'ex-libris de bibliothèque}ments which would\par
otherwise have a restricted reader\SGASevenSourceObstruction{l'ex-libris de bibliothèque}documented in the\par
scientific literature. -\par

\begin{center}
Author: LECTURES\par
Classmark: S/QA36.3\par
Book No.: 539502\par

UNIVERSITY OF NOTTINGHAM

George Green Library

of Science and Engineering
\end{center}

\noindent
Manuscripts\par
Manuscripts are reproduced by a photographic process; they must therefore be typed with extreme care.\par
Symbols not on the typewriter should be inserted by hand in indelible black ink. Corrections to the type-\par
script should be made by sticking the amended text over the old one, or by obliterating errors with white\par
correcting fluid. Authors receive 75 free copies.\par
The typescript is reduced slightly in size during reproduction; best results will not be obtained unless the\par
text on any one page is kept within the overall limit of 18 x 26.5 cm (7 x 10\ensuremath{\tfrac12} inches). The publishers will\par
be pleased to supply on request special stationery with the typing area outlined.\par
Manuscripts in English, German or French should be sent to Prof. Dr. A. Dold, Mathematisches Institut der\par
Universität Heidelberg, 69 Heidelberg/Germany, Tiergartenstraße or Prof. Dr. B. Eckmann, Eidgenössi-\par
sche Technische Hochschule, CH-8006 Zürich/Switzerland.\par

\medskip
\noindent
Lecture Notes in Physics\par
Bisher erschienen/Already published\par
Vol. 1: J. C. Erdmann, Wärmeleitung in Kristallen, theoretische Grundlagen und fortgeschrittene experimentelle Methoden. II, 283 Seiten. 1969. DM 20,--\par
Vol. 2: K. Hepp, Théorie de la renormalisation. III, 215 pages. 1969. DM 18,--\par
Vol. 3: A. Martin, Scattering Theory: Unitarity, Analyticity and Crossing. IV, 125 pages. 1969. DM 14,--\par
Vol. 4: G. Ludwig, Deutung des Begriffs »physikalische Theorie« und axiomatische Grundlegung der Hilbertraumstruktur der Quantenmechanik durch Hauptsätze des Messens. XI, 469 Seiten. 1970. DM 28,--\par
Vol. 5: M. Schaaf, The Reduction of the Product of Two Irreducible Unitary Representations of the Proper Orthochronous Quantummechanical Poincaré Group. IV, 120 pages. 1970. DM 14,--\par
Vol. 6: Group Representations in Mathematics and Physics. Edited by V. Bargmann. V, 340 pages. 1970. DM 24,--\par
Vol. 7: R. Balescu, J. L. Lebowitz, I. Prigogine, P. Résibois, Z. W. Salsburg, Lectures in Statistical Physics. V, 181 pages. 1971. DM 18,--\par
Vol. 8: Proceedings of the Second International Conference on Numerical Methods in Fluid Dynamics. Edited by M. Holt. IX, 462 pages. 1971. DM 28,--\par
Vol. 9: D. W. Robinson, The Thermodynamic Pressure in Quantum Statistical Mechanics. V, 115 pages. 1971. DM 14,--\par
Vol. 10: J. M. Stewart, Non-Equilibrium Relativistic Kinetic Theory. III, 113 pages. 1971. DM 14,--\par
Vol. 11: O. Steinmann, Perturbation Expansions in Axiomatic Field Theory. III, 126 pages. 1971. DM 14,--\par
\clearpage

% SGA7I-FRONT-P002 — controlling PDF page 3 — title page with library slips.
% Only unobscured visible fragments are transcribed; the hidden title is not reconstructed here.
\begin{center}
Lecture Notes in Mathematics

A collection of informal reports and seminars

Edited by A. Dold, Heidelberg and B. Eckmann, Zürich

\bigskip
288

\bigskip
Groupe\SGASevenSourceObstruction{les fiches de prêt de bibliothèque}

en Géor\SGASevenSourceObstruction{les fiches de prêt de bibliothèque}

\bigskip
Springer-Verlag

Berlin : Heidelberg : New York 1972
\end{center}
\clearpage

% SGA7I-FRONT-P003 — controlling PDF page 4 — copyright/title verso.
% Right-edge losses remain explicit; no standard-form notice is reconstructed.
\noindent
AMS Subject Classifications (1970): Primary: 14D05, 14D15, 14K99\par
\hspace*{8em}Secondary: 13D10, 14E20, 14F20\par

\medskip
\noindent
ISBN 3-540-05987-3 Springer-Verlag Berlin - Heidelberg - New York\par
ISBN 0-387-05987-3 Springer-Verlag New York - Heidelberg - Berlin\par

\medskip
\noindent
This work is subject to copyright. All rights are reserved, whether the whole or part of the material is concerne\SGASevenSourceGap{bord droit coupé}\par
specifically those of translation, reprinting, re-use of illustrations, broadcasting, reproduction by photocopying machin\SGASevenSourceGap{bord droit coupé}\par
or similar means, and storage in data banks.\par
Under § 54 of the German Copyright Law where copies are made for other than private use, a fee is payable to the publishe\SGASevenSourceGap{bord droit coupé}\par
the amount of the fee to be determined by agreement with the publisher.\par

\medskip
\noindent
© by Springer-Verlag Berlin : Heidelberg 1972. Library of Congress Catalog Card Number 72-88730. Printed in Germany.\par

\medskip
\noindent
Offsetdruck: Julius Beltz, Hemsbach/Bergstr.\par
\clearpage

% SGA7I-FRONT-P004 — controlling PDF page 5 — internal seminar title.
\begin{center}
SÉMINAIRE DE GÉOMÉTRIE ALGÉBRIQUE DU BOIS-MARIE

1967--1969

\bigskip
GROUPES DE MONODROMIE EN GÉOMÉTRIE ALGÉBRIQUE

(SGA 7 I)

\bigskip
Dirigé par A. GROTHENDIECK

avec la collaboration de

M. RAYNAUD et D. S. RIM
\end{center}
\clearpage

% SGA7I-FRONT-P005 — controlling PDF page 6 — Internet Archive provenance leaf.
\begin{center}
Digitized by the Internet Archive

in 2022 with funding from

Kahle/Austin Foundation

\bigskip
\texttt{https://archive.org/details/groupesdemonodro0288unse}
\end{center}
\clearpage

% SGA7I-FRONT-P006 — controlling PDF page 7 — Introduction, printed V inferred.
\begin{center}
INTRODUCTION
\end{center}

Soient $S$ un schéma et $f:X\longrightarrow S$ un morphisme de schémas. Si $f$ est
propre et lisse, et que le nombre premier $\ell$ est inversible sur $S$, les groupes de
cohomologie $\ell$-adique $H^i(X_{\bar s},\mathbb Z_\ell)$ des fibres géométriques de $f$ forment un système
local $\ell$-adique sur $S$. Pour $f$ seulement supposé propre, ces groupes sont les fibres
d'un faisceau $\ell$-adique $R^i f_*\mathbb Z_\ell$ sur $S$. La théorie des cycles évanescents met
en relation la ramification de ce faisceau sur $S$ et les singularités de $f$.

Nous ne considérerons que le cas où $S$ est un trait hensélien (= spectre d'un
anneau de valuation discrète hensélien). C'est en pratique le cas essentiel. Je renvoie
à (I 2.1) pour une description heuristique de la théorie. Soient $S=\operatorname{Spec}(V)$, $\bar\eta$
une clôture algébrique du corps des fractions $k(\eta)$ de $V$ et $\bar s$ la clôture algébrique
correspondante du corps résiduel $k(s)$ ($k(\bar s)$ est le corps résiduel du normalisé
$\bar V$ de $V$ dans $k(\bar\eta)$). Dans le cas particulier où $X$ est propre et plat sur
$S$, de dimension relative $n$, et où $f$ ne présente qu'un point de non lissité
$x\in X_s$ ; on définit des $\operatorname{Gal}(\bar\eta,\eta)$-modules $\phi^i$, de nature purement locale au voisinage
de $x$, nuls pour $i\notin[0,n]$ (I 4.2), et on construit une suite exacte longue
de $\operatorname{Gal}(\bar\eta,\eta)$-modules

\[
\cdots \longrightarrow H^i(X_{\bar s},\mathbb Z_\ell)
\xrightarrow{\mathrm{sp}} H^i(X_{\bar\eta},\mathbb Z_\ell)
\longrightarrow \phi^i
\longrightarrow H^{i+1}(X_{\bar s},\mathbb Z_\ell)
\longrightarrow \cdots
\]

(le $\operatorname{Gal}(\bar s,s)$-module $H^i(X_{\bar s},\mathbb Z_\ell)$ est regardé comme un
$\operatorname{Gal}(\bar\eta,\eta)$-module avec action triviale de l'inertie; $\mathrm{sp}$ est la flèche de spécialisation).

On donne aussi des critères, valables pour l'instant seulement en caractéristique
$0$, pour que les $\phi^i$, pour $i$ petits, soients nuls (par exemple, si $X$ est de plus
localement d'intersection complète, $\phi^i=0$ pour $i\ne n$ (I. 4.5)).
\clearpage

% SGA7I-FRONT-P007 — controlling PDF page 8 — Introduction, printed VI.
Le théorème de monodromie affirme que l'action du sous-groupe d'inertie $I$
de $\operatorname{Gal}(\bar\eta,\eta)$ est quasi-unipotente : pour $T\in I$, il existe des entiers $N>0$,
$M>0$ tels que l'endomorphisme $(T^M-I)^N$ de $H^i(X_{\bar\eta},\mathbb Z_\ell)$ soit nul. On donne de ce
théorème deux démonstrations. La première (I.1.2), par nature arithmétique, s'applique
dès que les groupes de cohomologie considérés ont des propriétés de finitude raisonnables.
Le point clef est que, lorsque $k(s)$ est de type fini sur son sous-corps
premier, l'action de $I$ est quasi-unipotente pour toute représentation $\ell$-adique de
$\operatorname{Gal}(\bar\eta,\eta)$. La seconde démonstration, plus géométrique, requiert la pureté et la résolution
des singularités : elle n'est pour l'instant valable qu'en caractéristique $0$
ou pour un $H^1$. Elle apporte de précieuses informations sur l'exposant de nilpotence
$N$ ($N\leq i+1$ pour un $H^i$). Ainsi qu'on le verra ultérieurement, elle se prête bien,
sur $\mathbb C$, à la comparaison avec la théorie transcendante.

Ces résultats, appliqués au $H^1$ des variétés abéliennes, i.e. à leur module
de Tate, permettent d'étudier la réduction mod $p$ de celles-ci. On donne ainsi deux
démonstrations du théorème de réduction stable des variétés abéliennes selon lequel,
après ramification, la fibre spéciale connexe du modèle de Néron est extension d'une
variété abélienne par un tore. La première (I. 6) reprend la méthode de la démonstration
arithmétique du théorème de monodromie. La seconde (IX 3.6) s'appuie sur une
analyse beaucoup plus fine du modèle de Néron, et de ses propriétés de polarisation.

Les exposés I à V de Grothendieck n'ont pas été rédigés. Ils ont été résumés
dans un exposé I. Les résultats énoncés y sont démontrés de façon succinte, mais
essentiellement complète.

L'exposé II applique la méthode des pinceaux de Lefschetz et (I 5.3) à
l'étude du groupe fondamental. Pour $S$ une surface sur un corps algébriquement clos
$k$, d'exposant caractéristique $p$, on montre que le groupe profini $\pi_1^{(p)}(S,s)$, complété
en dehors de $p$ du groupe fondamental, est de pro-$(p)$-présentation finie.

Comme expliqué plus haut, les exposés III à V n'existent pas.
\clearpage

% SGA7I-FRONT-P008 — controlling PDF page 9 — Introduction, printed VII.
L'exposé VI contient, avec quelques compléments, la théorie des déformations
de Schlessinger. On y prend soin de tenir compte des automorphismes infinitésimaux
des objets qu'on classifie,et de ne pas passer trop brutalement au foncteur des classes
d'isomorphie d'objets. On y donne aussi une nouvelle construction des
$\underline{\operatorname{Ext}}^{i}(L_{X/S},-)$ ($L_{X/S}$ complexe cotangent relatif de $X/S$). Cet exposé sert dans le
reste du séminaire surtout via l'étude qui y est faite des déformations des singularités
quadratiques ordinaires.

Les exposés VII et VIII sont consacrés à la théorie des biextensions. Cette
théorie joue un rôle essentiel dans l'exposé IX, pour exprimer ce qu'il advient d'une
polarisation quand on passe d'une variété abélienne à un modèle de Néron de celle-ci.
Cet exposé IX contient la démonstration du théorème de réduction stable des variétés
abéliennes, et diverses applications.

La suite de ce séminaire : SGA 7 II, par P. Deligne et N. Katz, paraîtra
ultérieurement.

\begin{flushright}
Bures sur Yvette, mai 1972,\\
P. DELIGNE
\end{flushright}
\clearpage

% SGA7I-FRONT-P009 — controlling PDF page 10 — blank verso.
% No intentional printed content. Faint reverse-side show-through only.
\null
\clearpage

% SGA7I-FRONT-P010 — controlling PDF page 11 — Table des Matières.
\begin{center}
Table des Matières
\end{center}

\noindent
Exposé I\par
\hspace*{2em}Résumé des premiers exposés de A. Grothendieck\par
\hspace*{2em}rédigé par P. Deligne \dotfill 1\par

\medskip
\noindent
Exposé II\par
\hspace*{2em}Propriétés de finitude du groupe fondamental\par
\hspace*{2em}par Michèle Raynaud \dotfill 25\par

\medskip
\noindent
Exposé VI\par
\hspace*{2em}Formal deformation theory\par
\hspace*{2em}by D.S. Rim \dotfill 32\par

\medskip
\noindent
Exposé VII\par
\hspace*{2em}Biextension de faisceaux de groupes\par
\hspace*{2em}par A. Grothendieck \dotfill 133\par

\medskip
\noindent
Exposé VIII\par
\hspace*{2em}Compléments sur les biextensions. Propriétés générales\par
\hspace*{2em}des biextensions des schémas en groupes\par
\hspace*{2em}par A. Grothendieck \dotfill 218\par

\medskip
\noindent
Exposé IX\par
\hspace*{2em}Modèles de Néron et monodromie\par
\hspace*{2em}par A. Grothendieck \dotfill 313\par
\clearpage

% SGA7I-FRONT-P011 — controlling PDF page 12 — blank verso.
% No intentional printed content. Faint reverse-side show-through only.
\null
\clearpage
